\begin{abstract}
Deploying specialized multi-task expert models on edge and IoT devices is heavily bottlenecked by severe storage and network transmission bandwidth constraints. While weight-space merging via Task Arithmetic enables multi-task capabilities by storing task-specific delta vectors, dense weights remain too large for resource-constrained wild environments. To address this bottleneck, we propose \textbf{Norm-Preserved Budgeted Task-Vector Pruning (NP-BTVP)}, a post-hoc weight sparsification and merging framework. We introduce norm-preserved deterministic pruning schemes: global \textbf{Uniform Pruning (NP-BTVP-U)} and layer-wise \textbf{Adaptive Saliency-Based Pruning (NP-BTVP-S)}. Crucially, we incorporate norm-preserving rescaling (scaling active updates by the reciprocal of their retention rate) as a deterministic signal-strength preservation heuristic, which mathematically prevents update norm shrinkage and preserves update signal strength. Empirically, we evaluate our methods across standard AdamW and Sharpness-Aware Minimization (SAM) optimizers. Surprisingly, we find that training-stage flatness (via SAM) does not inherently provide an additional coordinate-aligned pruning buffer compared to standard AdamW under well-converged regimes. However, under our norm-preserving rescaling, both standard AdamW and flatness-aware SAM experts demonstrate an extraordinary and nearly identical level of resilience to heavy sparsification, maintaining high classification accuracy under extreme sparsification budgets. Furthermore, under a tight 90\% sparsity budget, our deterministic Uniform Pruning (NP-BTVP-U) achieves \textbf{90.32\%} (SAM) and \textbf{90.34\%} (AdamW) average accuracy, performing remarkably close to the fully dense unpruned baseline (\textbf{91.00\%} and \textbf{90.94\%}) and competitively with the advanced stochastic DARE-Merging baseline, while outperforming TIES-Merging by \textbf{3.81\%} at half the parameter budget. Finally, we analyze the comparison between Uniform and Saliency pruning, demonstrating that layer-wise budget heterogeneity in NP-BTVP-S introduces scale instability during rescaling. This highlights the pragmatic superiority of the simpler, global NP-BTVP-U framework as a highly stable, robust, and cost-efficient solution for edge multi-task deployment.
\end{abstract}